
\section{Reproducibility}
\label{sec:repro}

The study was preregistered before P2 outcomes. The 180 formal samples, dual validators,
treatment contract, downstream policy, retrievers, top-$k$, hash dimension, RRF constant,
primary endpoint, McNemar test, bootstrap procedure, seed, and formal gates were frozen
before the confirmatory run.

P2 contains \PtwoEvaluations{} deterministic evaluations across six conditions and uses
zero provider/model calls and zero cash cost. The authoritative result artifact is
\texttt{research/evidence-surface-v1/p2/out/p2-results.json}, frozen at blob SHA
\texttt{e8a749a6707dbf4614b20ef6248f81eb43dd89d4}. The figure-value macros in this paper
record that exact source SHA.

The reported P2 results were produced by replaying the frozen GitHub artifacts and frozen
JavaScript logic in a V8 execution environment. A repository-local Node replay remains
desirable as a runtime-equivalence check. Such a replay cannot alter the frozen samples,
policy, retriever parameters, statistical rules, gates, or reported outcomes.

\section{Conclusion}
\label{sec:conclusion}

A downstream safety policy can only reason about the evidence it can see. In our formally
specified synthetic benchmark, hiding a required item raises unsupported decisive behavior
from 0\% to 16.67\%, but the effect is not generic: it is concentrated in
override/invalidation cases where the omitted superseding record leaves no visible trace
that anything is missing. Frozen top-$k$ retrieval produces UDRs near 30\%, and aggregate
mandatory-fact recall is not a strictly monotonic safety surrogate.

On an identical hybrid retrieval surface, a one-bit evidence-obligation completeness
signal removes the measured unsupported decisions while preserving decisive behavior
whenever formal coverage is complete. This does not imply that reliable coverage monitoring
is easy. It isolates the architectural requirement: systems that place safety checks
downstream of retrieval need a way to distinguish the visible evidence looking sufficient
from the required evidence obligations actually being covered.

The next research problem is therefore not only better retrieval. It is trustworthy
evidence-obligation accounting under open-world, changing state.
